
\begin{agradecimentos}

Agradeço primeiramente a Deus, por ter estado comigo desde o primeiro dia em que saí da casa dos meus pais. Sua presença foi o meu sustento nos momentos difíceis e em todo o processo de adaptação que um estudante que mora longe de casa enfrenta.
Aos meus pais, Marciano Savino e Mailza Savino, que, apesar de todos os problemas do cotidiano, nunca mediram esforços para me ajudar financeiramente e me oferecer suporte emocional. Se hoje cheguei onde estou, é porque pude me apoiar no ombro de gigantes que vocês são para mim.
À minha mulher, Thaiz Batista da Cunha, que sempre esteve comigo, fazendo-se presente mesmo nos momentos de distância física. Agradeço por ter mudado minha personalidade para melhor com seu jeito incrível de ser e por me ensinar a encontrar a felicidade nas pequenas coisas. Você fez toda a diferença nesta jornada.
Ao meu amigo Maycon Mendes Fernandes, pela parceria constante nas atividades e pela grande amizade que construímos. Maycon, você sempre foi minha primeira opção para os desafios acadêmicos, e sou grato por todo o conhecimento profissional e pessoal que adquiri ao seu lado.
Aos meus professores, em especial àqueles que tiveram paciência comigo nos primeiros anos de faculdade. O início foi um mundo totalmente novo, mas hoje me sinto integrado à área da tecnologia e devo grande parte dessa conquista aos ensinamentos e à dedicação de vocês.
Por fim, agradeço a mim mesmo, pelo esforço contínuo e pela resiliência que só eu conheço. Pelos dias intensos de estudo, pelas orações, pelas noites regadas a café e pelo desgaste mental superado. Olhando para trás, vejo que cada parte do processo valeu a pena para chegar até aqui.
\begin{flushright}
- Lucas Lorenzo Savino
\end{flushright}

Agradeço a Deus, por ter me dado a oportunidade de viver esta experiência acadêmica e por me guiar em cada passo do caminho. Sua presença constante foi fundamental para superar os desafios e alcançar este momento tão importante.
A minha mãe, Maria Cristina Mendes, que sempre me apoiou incondicionalmente, mesmo quando as circunstâncias eram difíceis. Sua força e amor foram essenciais para que eu pudesse seguir em frente, mesmo nos momentos mais desafiadores.
A minha namorada, Anna Carolina F. Sartori, que esteve ao meu lado durante toda a jornada, oferecendo apoio emocional e compreensão a todo momento. Sua presença foi um porto seguro nos momentos de estresse e uma fonte de motivação para continuar avançando com leveza.
Ao meu amigo Lucas Lorenzo Savino, com quem compartilhei muitos momentos de estudo e aprendizado. A parceria e a amizade que construímos foram fundamentais para o meu desenvolvimento acadêmico e pessoal. Lucas, agradeço por todo o conhecimento que adquiri com você em nossas discussões.
Aos meus amigos e colegas de curso, que tornaram a experiência universitária mais leve e divertida. Em especial, a Rodrigo P. Almeida, Lucas R. Rangel, Pedro M. Menezes e Ester P. Nunes, vocês foram uma parte importante da minha jornada, e sou grato por cada momento compartilhado.
Por fim, aos meus professores, que me ensinaram não apenas o conteúdo acadêmico, mas também lições de vida e ética profissional. Agradeço por todo o conhecimento transmitido e pela dedicação em nos preparar para o futuro.
\begin{flushright}
- Maycon Mendes Fernandes
\end{flushright}

\end{agradecimentos}


